\documentclass[11pt]{article}

\usepackage[T1]{fontenc}
\usepackage{lmodern}
\usepackage[margin=1in]{geometry}
\usepackage{amsmath,amssymb,mathtools,amsthm}
\usepackage{microtype}
\usepackage[hidelinks]{hyperref}
\usepackage{url}

\allowdisplaybreaks[2]

\newtheorem{theorem}{Theorem}[section]
\newtheorem{proposition}[theorem]{Proposition}
\newtheorem{lemma}[theorem]{Lemma}
\newtheorem{corollary}[theorem]{Corollary}

\DeclareMathOperator{\tr}{tr}
\newcommand{\F}{\mathrm F}
\newcommand{\R}{\mathbb R}
\newcommand{\cS}{\mathcal S}
\newcommand{\cP}{\mathcal P}
\newcommand{\cD}{\mathcal D}

\title{More than 67.302\% of the zeros of the Riemann zeta function are simple and lie on the critical line}
\date{}

\begin{document}

\maketitle

\begin{abstract}
Let $N(T,2T)$ count the zeros of the Riemann zeta function with ordinates in
$(T,2T]$, with multiplicity, and let $N_0^s(T,2T)$ count the simple zeros on
the critical line in that interval.  We refine the rank--trace argument of
Claude~\cite{Claude26} and prove
\[
 \liminf_{T\to\infty}\frac{N_0^s(T,2T)}{N(T,2T)}
 \ge 0.673025467453\ldots .
\]
The proof uses the analytic estimates of Theorem~D in~\cite{Claude26} and a
rigorously verified seven-point inequality for the simple-zero vectors.
\end{abstract}

\section{Statement and setup}

We use the notation and normalisations of~\cite[\S\S2,4,7.1]{Claude26}.  Put
\[
 \ell:=\log\frac{T}{2\pi},\qquad I:=(T,2T],\qquad
 D_0:=T^{1/2},\qquad I':=(T-D_0,2T+D_0].
\]
For the optimised test family of Theorem~D, $L=\ell$ and
\[
 \phi(u)=\cos\!\left(\frac{\sqrt2\,u}{\ell}\right)^{1/2}
 \varrho(L/2-|u|),
\]
where the fixed-width ramp $\varrho$ is defined in~\cite[\S2.2]{Claude26}.
Let
\[
 a:=\frac1L\int_{\R}\phi(u)^2\,du,
 \qquad \widehat G:=\frac{G}{aL^2},
 \qquad \widehat A:=\frac{A}{aL^2},
 \qquad d:=\left\lfloor\frac{LT}{2\pi}\right\rfloor.
\]
As in~\cite[\S4.1]{Claude26}, divide the zeros in $I'$ into the simple
critical-line zeros $\cS_1$, the distinct multiple critical-line zeros
$\cS_2$, and the unordered off-line pairs $\cP$.  Write
\[
 s_1:=\#\cS_1,\qquad s_2:=\#\cS_2,\qquad p:=\#\cP.
\]
Then
\begin{equation}\label{eq:Nlower}
 N(I')\ge s_1+2s_2+2p.
\end{equation}
For $\rho\in\cS_1$, define
\[
 u_\rho:=\bigl(\widehat\phi(\gamma_\rho-\tau_k)\bigr)_{0\le k<d}
 \in\R^d,
 \qquad
 v_\rho:=\frac{u_\rho}{\sqrt a\,L}.
\]
The simple-zero part of $\widehat A$ is
\[
 P_1=VV^{\mathsf T}=\sum_{\rho\in\cS_1}v_\rho v_\rho^{\mathsf T},
\]
where the columns of $V$ are the $v_\rho$.  By Lemma~2.2
of~\cite{Claude26}, $\|v_\rho\|\le1$.  For
\[
 Q':=\widehat A-P_1,
\]
the decomposition in the proof of Proposition~4.4 of~\cite{Claude26} gives
\begin{equation}\label{eq:index}
 n_+(Q')\le s_2+p.
\end{equation}

Set
\[
 H_{\mathrm{MT}}
 :=\frac32-\frac1{\sqrt2}\cot\frac1{\sqrt2}
 =0.6725007036794116457\ldots .
\]
Theorem~D of~\cite{Claude26} proves the lower bound $H_{\mathrm{MT}}$.
We prove the following refinement.

\begin{theorem}\label{thm:main}
One has
\[
 \liminf_{T\to\infty}
 \frac{N_0^s(T,2T)}{N(T,2T)}
 \ge
 \frac{267\cdot10^{9}H_{\mathrm{MT}}-5.32\cdot10^{8}}{266{,}001{,}357{,}363}
 =0.6730254674531403414\ldots .
\]
\end{theorem}

\section{A stability refinement of the rank--trace inequality}

Define the convex function
\begin{equation}\label{eq:Psi}
 \Psi(t):=
 \begin{cases}
  (t-1)^2,&0\le t\le2,\\
  2t-3,&t\ge2.
 \end{cases}
\end{equation}
For a positive semidefinite matrix $M$, let
\[
 \cD(M):=\tr\Psi(M).
\]

\begin{lemma}[Stability-enhanced rank--trace inequality]\label{lem:stable-rank-trace}
Let $V$ be a $d\times r$ matrix whose columns have norm at most one.  Put
$P=VV^*\succeq0$ and $M=V^*V$, and let $Q$ be Hermitian with $n_+(Q)\le b$.
Then
\begin{equation}\label{eq:stable-rank-trace}
 \|P+Q\|_{\F}^2
 \ge 4\tr(P+Q)-3r-4b+\cD(M).
\end{equation}
\end{lemma}

\begin{proof}
Write $Q=Q_+-Q_-$ with $Q_\pm\succeq0$, $Q_+Q_-=0$, and
$\operatorname{rank}Q_+\le b$.  Since $P,Q_+\succeq0$,
\[
 \|P+Q\|_{\F}^2
 \ge \|P-Q_-\|_{\F}^2+\|Q_+\|_{\F}^2.
\]
If $q_j>0$ are the nonzero eigenvalues of $Q_+$, then
$q_j^2\ge4q_j-4$, so
\begin{equation}\label{eq:qplus}
 \|Q_+\|_{\F}^2\ge4\tr Q_+-4b.
\end{equation}
Let $p_1\ge\cdots\ge p_r\ge0$ be the eigenvalues of $M$ (the nonzero
eigenvalues of $P$, padded by zeros), and let $n_1\ge n_2\ge\cdots\ge0$
be the eigenvalues of $Q_-$.  Von Neumann's trace inequality gives
$\tr(PQ_-)\le\sum_{i\le r}p_in_i$, hence
\[
 \|P-Q_-\|_{\F}^2+4\tr Q_-
 \ge \sum_{i\le r}\bigl((p_i-n_i)^2+4n_i\bigr).
\]
For $p\ge0$,
\[
 \min_{n\ge0}\bigl((p-n)^2+4n\bigr)
 =\begin{cases}p^2,&0\le p\le2,\\4p-4,&p\ge2,\end{cases}
 =2p-1+\Psi(p).
\]
Therefore
\begin{equation}\label{eq:qminus}
 \|P-Q_-\|_{\F}^2
 \ge2\tr P-r+\cD(M)-4\tr Q_-.
\end{equation}
Combining \eqref{eq:qplus} and \eqref{eq:qminus} gives
\[
 \|P+Q\|_{\F}^2
 \ge4\tr(P+Q)-2\tr P-r-4b+\cD(M).
\]
Since $\tr P$ is the sum of the squared column norms of $V$, $\tr P\le r$.
\end{proof}

\begin{corollary}\label{cor:count}
With $M:=V^{\mathsf T}V$,
\begin{equation}\label{eq:s1}
 s_1\ge4\tr\widehat A-\|\widehat A\|_{\F}^2-2N(I')+\cD(M).
\end{equation}
After removing the tails as in Proposition~4.2 of~\cite{Claude26},
\begin{equation}\label{eq:global-defect}
 N_0^s(T,2T)
 \ge H_{\mathrm{MT}}N(T,2T)+\cD(M^\circ)-o(N(T,2T)),
\end{equation}
where $M^\circ$ is the Gram matrix of the retained central simple zeros.
\end{corollary}

\begin{proof}
Apply Lemma~\ref{lem:stable-rank-trace} to
$P_1+Q'=\widehat A$.  Equations \eqref{eq:Nlower} and \eqref{eq:index}
give
\[
 \|\widehat A\|_{\F}^2
 \ge4\tr\widehat A-3s_1-4s_2-4p+\cD(M)
 \ge4\tr\widehat A-s_1-2N(I')+\cD(M),
\]
which is \eqref{eq:s1}.

Use the same comparison $\widehat A=\widehat G-\widehat E$ and tail
estimates as in Propositions~4.2 and~4.4 of~\cite{Claude26}.  Pinching the
full Gram matrix into the retained central block and its complement gives
$\cD(M)\ge\cD(M^\circ)$ because $X\mapsto\tr\Psi(X)$ is convex and
unitarily invariant.  Moreover $N(I')=N(T,2T)+o(N(T,2T))$, and Theorem~D
gives
\[
 \tr\widehat G=N(T,2T)(1+o(1)),\qquad
 \|\widehat G\|_{\F}^2
 =\left(\frac1{c_1^*}+o(1)\right)N(T,2T),
\]
where $2-1/c_1^*=H_{\mathrm{MT}}$.  Substitution in \eqref{eq:s1}
yields \eqref{eq:global-defect}.
\end{proof}

\section{The Montgomery--Taylor overlap kernel}

Put $h:=2\pi/L$ and define the normalised ordinate
\[
 x_\rho:=\frac{\gamma_\rho-T}{h}
 =\frac{L(\gamma_\rho-T)}{2\pi}.
\]

\begin{lemma}\label{lem:kernel-limit}
Fix $R_0>0$.  Delete the simple zeros in strips of normalised width $L^2$
at the two ends of the grid.  The number deleted is $o(N(T,2T))$, and
uniformly for retained simple zeros with $|x_\rho-x_{\rho'}|\le R_0$,
\begin{equation}\label{eq:kernel-limit}
 \langle v_\rho,v_{\rho'}\rangle
 =k(x_\rho-x_{\rho'})+o(1),
\end{equation}
where
\begin{align}
 k(x)&:=\frac{K(x)}{K(0)},\label{eq:kdef}\\
 K(x)&:=\int_{-1/2}^{1/2}\cos(\sqrt2\,t)\cos(2\pi xt)\,dt\notag\\
 &=\frac{\sin(\pi x-1/\sqrt2)}{2\pi x-\sqrt2}
   +\frac{\sin(\pi x+1/\sqrt2)}{2\pi x+\sqrt2},\label{eq:Kexplicit}
\end{align}
and $K(0)=\sqrt2\sin(1/\sqrt2)$.  The singularities in
\eqref{eq:Kexplicit} are removable.
\end{lemma}

\begin{proof}
For the full Gabor grid, Lemma~2.2 of~\cite{Claude26} gives
\[
 \frac1{aL^2}\sum_{j\in\mathbb Z}
 \widehat\phi(\gamma-\tau_j)\widehat\phi(\gamma'-\tau_j)
 =\frac{\Phi(\gamma-\gamma')}{aL},
 \qquad \Phi=\widehat{\phi^2}.
\]
For points at normalised distance at least $L^2$ from the grid ends, the
terms with $j\notin[0,d)$ contribute $O(L^{-2})$, uniformly when
$|x_\rho-x_{\rho'}|$ is bounded, by the $r^{-2}$ decay used
in~\cite[\S5.3]{Claude26}.

For fixed $x$, substitution $u=Lt$ gives the normalised full-grid overlap
\begin{equation}\label{eq:normalised-overlap}
 \frac{\Phi(hx)}{aL}
 =
 \frac{\displaystyle\int_{-1/2}^{1/2}
       \phi(Lt)^2 e^{-2\pi ixt}\,dt}
      {\displaystyle\int_{-1/2}^{1/2}\phi(Lt)^2\,dt}
 \longrightarrow \frac{K(x)}{K(0)}.
\end{equation}
Indeed,
$\phi(Lt)^2\to\cos(\sqrt2\,t)\mathbf1_{[-1/2,1/2]}(t)$ in $L^1$;
the fixed-width taper has rescaled length $O(L^{-1})$.  The convergence in
\eqref{eq:normalised-overlap} is uniform for $x$ in compact sets.  Together
with the tail estimate, this proves \eqref{eq:kernel-limit}.  Finally, the
deleted strips have ordinate length $O(L)$ and contain $O(L^2)=o(N(T,2T))$
zeros, since $N(t+1)-N(t)\ll\log(t+3)$.
\end{proof}

Set
\[
 w(x):=k(x)^2.
\]

\section{A seven-point local inequality}

For nonnegative gaps $g_1,\ldots,g_6$, define
\begin{equation}\label{eq:F6}
 \mathcal F_6(g_1,\ldots,g_6)
 :=\frac1{3000}\sum_{i=1}^6g_i
 +\sum_{r=1}^6\frac{2}{7-r}
   \sum_{i=1}^{7-r}w(g_i+\cdots+g_{i+r-1}).
\end{equation}
The second term contains all $21$ pairwise separations among seven ordered
points.

\begin{proposition}[Certified local inequality]\label{prop:F6}
For every $g_1,\ldots,g_6\ge0$,
\begin{equation}\label{eq:F6bound}
 \mathcal F_6(g_1,\ldots,g_6)\ge\beta_0:=\frac{3{,}826{,}217}{10^{9}}.
\end{equation}
\end{proposition}

\begin{proof}[Verification]
The verifier at
\begin{center}
 \url{https://github.com/npip99/zeta-zeros}
\end{center}
uses Arb interval arithmetic through \texttt{python-flint}.  It evaluates
the entire-sinc expression
\[
 K(x)=\frac12\left(
 \frac{\sin(\pi x-1/\sqrt2)}{\pi x-1/\sqrt2}
 +\frac{\sin(\pi x+1/\sqrt2)}{\pi x+1/\sqrt2}
 \right)
\]
at 128-bit precision and constructs rigorous lower bounds for $w$ on a grid
of mesh $1/8000$.  If $\sum g_i\ge3000\beta_0=11.478651$, the linear term
proves the claim.
On the remaining compact region, the one-variable term
\[
 U(g):=\frac g{3000}+\frac13w(g)
\]
reduces each coordinate to three certified unions of grid cells.  The
resulting $3^6$ boxes are proved by interval lower bounds, a certified convex
tangent bound when applicable, or bisection of the widest coordinate.  The
program fails if a terminal box remains unresolved.

The committed certificate records \texttt{verified=true}, grid $8000$,
precision $128$ bits, $963{,}245$ visited nodes, and maximum depth $69$.
The adjacent target $3{,}826{,}218/10^{9}$ fails at a terminal cell, so the
certified target is maximal on the $10^{-9}$ lattice at this grid.
It also records the hashes of the interval tables used in the search.
\end{proof}

\begin{lemma}\label{lem:block-energy}
Let $y_1<\cdots<y_m$ with $m\ge7$, and put
\[
 E_m:=2\sum_{1\le i<j\le m}w(y_j-y_i).
\]
Then
\begin{equation}\label{eq:block-energy}
 E_m+\frac1{500}(y_m-y_1)
 \ge\beta_0(m-6).
\end{equation}
\end{lemma}

\begin{proof}
Write $g_i=y_{i+1}-y_i$ and sum \eqref{eq:F6bound} over the $m-6$
consecutive seven-point windows.  A pair spanning $r$ gaps occurs at most
$7-r$ times and has coefficient $2/(7-r)$ each time, so its total
contribution is at most $2w(y_j-y_i)$.  Each single gap occurs at most six
times, giving at most
\[
 \frac6{3000}\sum_{i=1}^{m-1}g_i
 =\frac1{500}(y_m-y_1).
\]
\end{proof}

\begin{lemma}\label{lem:block-defect}
If $G\succeq0$ is Hermitian, then
\begin{equation}\label{eq:block-defect}
 \tr\Psi(G)
 \ge\min\left\{1,\,2\sum_{i<j}|G_{ij}|^2\right\}.
\end{equation}
\end{lemma}

\begin{proof}
If every eigenvalue of $G$ is at most $2$, then
$\Psi(G)=(G-I)^2$, and
\[
 \tr\Psi(G)=\tr(G-I)^2\ge2\sum_{i<j}|G_{ij}|^2.
\]
If some eigenvalue $\lambda>2$, then $\Psi(\lambda)=2\lambda-3>1$.
\end{proof}

Take
\begin{equation}\label{eq:mA}
 m:=267,
 \qquad
 A_0:=\beta_0(m-6)=\frac{998{,}642{,}637}{10^{9}}<1.
\end{equation}
Let $B$ be a block of $267$ consecutive retained simple zeros.  Write $G_B$
for its Gram matrix.  If $\operatorname{span}(B)/500\ge A_0$, then
\eqref{eq:269block} below is immediate.  Otherwise
$\operatorname{span}(B)<500$, so Lemma~\ref{lem:kernel-limit} applies
uniformly to every pair in $B$.  Since $m$ is fixed,
\[
 2\sum_{i<j}|(G_B)_{ij}|^2=E_m+o(1).
\]
Lemmas~\ref{lem:block-energy} and~\ref{lem:block-defect} give, uniformly in
$B$,
\begin{equation}\label{eq:269block}
 \cD(G_B)+\frac1{500}\operatorname{span}(B)
 \ge A_0-o(1).
\end{equation}

\section{Shifted block pinching}

Let
\[
 x_1<\cdots<x_{S^\circ}
\]
be the normalised ordinates of the retained central simple zeros.  Then
\begin{equation}\label{eq:Scentral}
 S^\circ=N_0^s(T,2T)-o(N(T,2T)).
\end{equation}
For each of the $m=267$ offsets, partition the ordered points into full
consecutive blocks of size $m$, leaving at most $2(m-1)$ endpoint points.
Pinching $M^\circ$ to the full principal blocks and the remaining block gives
\begin{equation}\label{eq:pinching}
 \cD(M^\circ)\ge\sum_B\cD(G_B).
\end{equation}
This follows because pinching is an average of unitary conjugations and
$X\mapsto\tr\Psi(X)$ is convex and unitarily invariant.

For each offset, sum \eqref{eq:269block} over the full blocks and use
\eqref{eq:pinching}; then average over the $m$ offsets.  The average number of
full blocks is $S^\circ/m+O(1)$.  Each gap $x_{j+1}-x_j$ belongs to a block
span for at most $m-1$ offsets, and
\[
 x_{S^\circ}-x_1
 \le\frac{T}{h}
 =\frac{LT}{2\pi}
 =d+O(1)
 =N(T,2T)+o(N(T,2T))
\]
by the Riemann--von Mangoldt formula.  The blockwise $o(1)$ term in
\eqref{eq:269block} is uniform, so its total is $o(N(T,2T))$.  Therefore
\begin{equation}\label{eq:defect-global}
 \cD(M^\circ)
 \ge \frac{A_0}{m}N_0^s(T,2T)
 -\frac{m-1}{500m}N(T,2T)-o(N(T,2T)).
\end{equation}
With $m=267$ and $A_0=998{,}642{,}637/10^{9}$,
\begin{equation}\label{eq:defect-numbers}
 \cD(M^\circ)
 \ge
 \frac{332{,}880{,}879}{89\cdot10^{9}}N_0^s(T,2T)
 -\frac{133}{66{,}750}N(T,2T)-o(N(T,2T)).
\end{equation}

Substitute \eqref{eq:defect-numbers} into \eqref{eq:global-defect}.  This gives
\begin{align*}
 N_0^s(T,2T)
 &\ge H_{\mathrm{MT}}N(T,2T)
 +\frac{332{,}880{,}879}{89\cdot10^{9}}N_0^s(T,2T)\\
 &\qquad-\frac{133}{66{,}750}N(T,2T)-o(N(T,2T)).
\end{align*}
Hence
\[
 \left(1-\frac{332{,}880{,}879}{89\cdot10^{9}}\right)N_0^s(T,2T)
 \ge
 \left(H_{\mathrm{MT}}-\frac{133}{66{,}750}\right)N(T,2T)
 -o(N(T,2T)).
\]
Divide by $N(T,2T)$ and let $T\to\infty$ to obtain
\[
 \liminf_{T\to\infty}\frac{N_0^s(T,2T)}{N(T,2T)}
 \ge
 \frac{267\cdot10^{9}H_{\mathrm{MT}}-5.32\cdot10^{8}}{266{,}001{,}357{,}363}
 =0.6730254674531403414\ldots .
\]
This proves Theorem~\ref{thm:main}.

\section*{Reproducibility}

Proposition~\ref{prop:F6} is the only computer-assisted input.  Its source,
tests, trust-base description, and certificate are available at
\begin{center}
\url{https://github.com/npip99/zeta-zeros}.
\end{center}
The verifier reconstructs every transcendental interval enclosure from the
stated formulas.

\begin{thebibliography}{99}

\bibitem{Claude26}
Claude,
\emph{More than two thirds of the zeros of the Riemann zeta function lie on the critical line},
preprint, August 10, 2026,
\href{https://www-cdn.anthropic.com/564f962e60643842f5fcb4a17c9dbc8f608f1c37.pdf}
{available online}.

\bibitem{Ainta26}
ainta,
\emph{zeta-simple-zeros: a 67.30085\% lower bound for simple zeros of the Riemann zeta function},
research artifact, August 10, 2026,
\href{https://github.com/ainta/zeta-simple-zeros}{available online}.
The argument and verifier of the present note originate there; only the
certified target in Proposition~4.1 and the propagated constants differ.

\bibitem{Johansson17}
F.~Johansson,
\emph{Arb: efficient arbitrary-precision midpoint-radius interval arithmetic},
IEEE Trans. Comput. \textbf{66} (2017), no.~8, 1281--1292.

\bibitem{Mon73}
H.~L.~Montgomery,
\emph{The pair correlation of zeros of the zeta function},
in: Analytic Number Theory (St. Louis, 1972),
Proc. Sympos. Pure Math. \textbf{24}, Amer. Math. Soc., 1973, 181--193.

\bibitem{Mon75}
H.~L.~Montgomery,
\emph{Distribution of the zeros of the Riemann zeta function},
in: Proceedings of the International Congress of Mathematicians
(Vancouver, 1974), Vol.~1,
Canad. Math. Congress, 1975, 379--381.

\end{thebibliography}

\end{document}
